\documentclass{beamer}
\usetheme{Berlin}
\usecolortheme{wolverine}
\usepackage[spanish]{babel}
\usepackage{graphicx}
\usepackage{amsmath}

\title{Visualización y análisis de funciones de dos variables: dominio, curvas de nivel, trazas y tasas de cambio usando MATLAB}
\author{Nora Carolina Servellón}
\institute[EMCC-UNAH]{\textit{Escuela de Matem\'aticas y Ciencias de la Computaci\'on}\\ \textbf{Universidad Aut\'onoma de Honduras}}
\date{Abril 2026}
\logo{\includegraphics[width=0.2\textwidth]{img/unah.png}
\includegraphics[width=0.12\textwidth]{img/facultad.jpeg}}

\begin{document}

% 1: Portada
\begin{frame}
\titlepage
\end{frame}

% 2: Índice 
\begin{frame}{Contenido e Introducción}
\tableofcontents
\end{frame}

% 3: Introducción
\begin{frame}{Introducción}
\textbf{Introducción:} Las funciones de dos variables \(z=f(x,y)\) son fundamentales en cálculo multivariable. Sin embargo, visualizar su dominio, forma y comportamiento es un reto. Con MATLAB exploraremos gráficas 3D, curvas de nivel y tasas de cambio para comprender conceptos como pendiente y dirección de máximo ascenso.
\begin{figure}
    \centering
    \includegraphics[width=0.25\linewidth]{img/matlab.png}
\end{figure}
\end{frame}


% 4: Conceptos clave
\section{Conceptos clave}
\begin{frame}{Conceptos clave}
\begin{itemize}
\item \textbf{Dominio}: región del plano donde la función está definida.
\item \textbf{Gráfica}: superficie \(z=f(x,y)\) en el espacio 3D.
\item \textbf{Trazas verticales}: intersección con planos \(x=a\) o \(y=b\).
\item \textbf{Curvas de nivel}: conjunto \(f(x,y)=c\) en el plano \(xy\).
\item \textbf{Tasa de cambio}: variación de \(f\) al moverse en una dirección.
\end{itemize}
\end{frame}


% 5: Dominio y superficies
\section{Domionio y Superficies}
\begin{frame}{Dominio y superficies}
\begin{columns}[T]
\column{0.33\textwidth}
\textbf{Dominio:} 
\vspace{1cm}
\includegraphics[width=1\textwidth]{img/dominio.eps}
\column{0.33\textwidth}
\textbf{Superficies:}\\
\includegraphics[width=1\textwidth]{img/paraboloide.eps}\\
Paraboloide elíptico \(2x^2+5y^2\)\\
\column{0.33\textwidth}
\includegraphics[width=1\textwidth]{img/hiperbolico.eps}\\
Paraboloide hiperbólico\\ \(x^2-3y^2\)
\end{columns}
\end{frame}

% 6: Curvas de nivel y pendiente
\section{Curvas de nivel y pendiente}
\begin{frame}{Curvas de nivel y pendiente}
\begin{columns}[T]
\column{0.35\textwidth}
\textbf{Mapas de contorno:}
\includegraphics[width=0.9\textwidth]{img/contorno_parab.eps}\\
Curvas de nivel del paraboloide elíptico
\column{0.35\textwidth}
\vspace{0.4cm}
\includegraphics[width=0.9\textwidth]{img/contorno_hip.eps}\\
Curvas de nivel del paraboloide hiperbólico
\vspace{0.5cm}
\column{0.3\textwidth}
\textbf{Relación:} A menor espaciamiento entre curvas, mayor pendiente de la superficie.
\end{columns}
\end{frame}

% 7: Tasa de cambio y gradiente
\section{Tasa de cambio y gradiante}
\begin{frame}{Tasa de cambio y gradiente}
\begin{itemize}
\item Sobre el plano \(f(x,y)=12-2x-3y\):
\begin{align*}
\text{Tasa}(A\to B) &= -2,\quad A=(0,0), B=(1,0)\\
\text{Tasa}(A\to C) &= -3,\quad C=(0,1)
\end{align*}
\item \textbf{Gradiente numérico} en \(A\): \(\nabla f \approx (-2,-3)\)
\item \textbf{Dirección de máximo ascenso}: \(\frac{(2,3)}{\sqrt{13}}\)
\end{itemize}
\begin{center}
\includegraphics[width=0.35\textwidth]{img/ascenso.eps}
\end{center}
\end{frame}

% 8: Conclusiones 
\section{Conclusiones}
\begin{frame}{Conclusiones}
\begin{itemize}
\item MATLAB permite visualizar dominio, superficies y curvas de nivel eficazmente.
\item El espaciamiento de las curvas de nivel indica la pendiente de la superficie.
\item La tasa de cambio depende de la dirección; el gradiente señala el máximo ascenso.
\item El código desarrollado es reusable y educativo.
\end{itemize}
\end{frame}

% 8: Conclusiones 
\begin{frame}{Bibliografía}
\section{Bibliograía}
\begin{itemize}
\item J. Stewart, \emph{Cálculo de varias variables}, Cengage, 2012.
\item MathWorks, \emph{MATLAB Documentation}.
\item Código fuente: \url{https://github.com/servelloncarol/proped-3}
\end{itemize}
\end{frame}

\end{document}
